\documentclass[12pt,letterpaper]{article}
\usepackage[utf8]{inputenc}
\usepackage[T1]{fontenc}
\usepackage{geometry}
\usepackage{setspace}
\usepackage{titlesec}
\usepackage{fancyhdr}
\usepackage{enumitem}
\usepackage{parskip}
\usepackage{hyperref}

% Page setup
\geometry{margin=1in}
\setlength{\parindent}{0pt}
\setlength{\parskip}{6pt}

% Header and footer
\pagestyle{fancy}
\fancyhf{}
\fancyhead[L]{\textsc{Markets and Bubbles}}
\fancyhead[R]{\thepage}
\renewcommand{\headrulewidth}{0.4pt}

% Title formatting
\titleformat{\section}{\Large\bfseries}{\thesection}{1em}{}
\titleformat{\subsection}{\large\bfseries}{\thesubsection}{1em}{}

% Document info
\title{\Huge\textbf{Markets and Bubbles}}
\author{Mahir Bansal}
\date{}

\begin{document}

\maketitle

\vspace{2cm}

\section*{Introduction}

I pay twenty dollars a month for Cursor. Cursor pays OpenAI/Claude/xAI for API calls. Those calls run on GPUs rented from a cloud that bought them by the truckload. Nvidia pays TSMC to fabricate them. TSMC pays ASML for the machines that etch circuits smaller than dust. ASML pays the firms that ship tin and argon across oceans.

Every prompt stretches down the entire value chain, from a keyboard to a lithography machine in a clean room.

Everyone in the chain sees margin somewhere else. Cursor looks to the model provider. The model provider points to the cloud. The cloud points to Nvidia. Nvidia points to TSMC. TSMC points to ASML. At the bottom, materials companies move tin and argon across oceans.

Between those nodes sit a lot of hands. Training data providers. Data center builders. Power plants. Grid operators. Undersea cables. Cooling vendors. Real estate trusts. EDA software. Wafer suppliers. There are more players than any speaker series or conference can hold.

Follow the money long enough and you end up staring at a factory, a power station, and a balance sheet that doesn't care about hype. That's why bubbles can be clarifying. They force an accounting. ``Who's actually making money?''

\section{Where the margin lands}

Margin tends to stick at the nodes you can't go around.

\begin{enumerate}[label=\textbf{\arabic*.}]
    \item \textbf{Deployments}\\
    Products with direct relationships to users. Most people haven't heard of AWS, they have heard of Netflix. These are firms like Cursor. TurboScribe. Notion. ChatGPT. Grok. A few make themselves unavoidable enough to hold margin. Most leak value to the next node betting on future dollars for market share today.

    \item \textbf{Platforms}\\
    Model providers turn capital expense into API endpoints and price the difference. Their edge is R\&D talent to fuel increasing model capabilities. Their risk is commoditization. Their defense is performance per dollar and an ecosystem that makes them unavoidable. Where leaving feels like loss.

    \item \textbf{Picks and shovels}\\
    Chips. Fabs. Tools. Materials. Energy. In the gold rush, everyone moved around different mines, but irrespective of the mine they were going to they all needed picks, shovels, and jeans. That was Levi's. These scale with ambition but don't require picking a single winner. Only that digging continues tomorrow.

    \item \textbf{Sovereigns}\\
    States bankrolling the chain with grants, loans, and credits. They bet that winning the race is important, and losing money today is worth limitless opportunity and not losing tomorrow.
\end{enumerate}

The stack taxes you at each layer. To keep margin, players become unavoidable by physics, switching costs, or taste.

\section{Experiments}

Bubbles are experiments.

We call them irrational because our spreadsheet lives in the present. But markets don't price the present. They compute possibilities of different futures and pay to run them faster than history would have allowed.

Think of this as a voting system. Everyone is betting on their vision of tomorrow, backing it up with dollars today. Something feels like a bubble when this system becomes sure that something relatively minor today will prove major tomorrow. That doesn't mean they're all hype. Every node is a ``what if this is the future?''. Most nodes crash but the survivors may become the bedrock for everything.

\section{Past experiments}

\textbf{Tulips} looked absurd. Flowers priced above houses. Underneath was a test of whether rarity and status could be financialized. The instruments failed. A few ideas stayed. Futures learned a couple tricks.

\textbf{Railroads} ruined many investors and gave a continent a circulatory system. The tracks rewired the map. Return to society can be a dividend equity never sees.

\textbf{Dot com} made attention an asset and bandwidth a right. Most sites died. Some did not. Broadband and data centers. Ecommerce platforms.

\textbf{Crypto}'s experiment is still very much ongoing. As prices fluctuate, we're testing futures with decentralized systems. Many nodes will die, some will survive. I think the schema will survive.

It's easy to talk about these things. Run countless balls through a Galton Board, we'll know the relative distribution of where they'll fall as a group. Many will die, some will survive. Wait a while, and we'll know where each landed. It's a lot harder to make an educated guess of where each ball will land while the experiment is still running.

\section{The AI experiment}

This experiment is big. The value chain stretches across all industries at all stages. We're testing visions of a future where cognition isn't unique to us. I believe that has to be the case and it looks like the market does too. The question is what that future looks like. We can't know that.

There are some things I think we can know, though:

\begin{itemize}
    \item Energy. Without cheap reliable power there is no thinking at scale.
    \item Compute is leverage. Someone needs to convert joules into tokens efficiently.
    \item Models. Some will be general, some specialist.
    \item Memory. Context. The best model can only compute as much input as it's given. Cars don't just need incredible engines, they need gas and a tank.
    \item Workflows. People use cars not solo tires. Someone needs to deploy all of the stack into products and services that feel like magic. Unavoidable ones.
\end{itemize}

The value chain needs all these nodes and the ones in between them I didn't mention.

\section{Governments}

The AI race is a market story and a state story.

Governments support fabs, transmission, and foundational labs because model capabilities translate to national capabilities. Depend on another country for thinking and you let its government control everything you touch.

Maybe this future we're testing is so important, it'll mean changing the economy to make the economics work. Maybe it's too important to fail.

\section{Final Thoughts}

We're in the middle of one of the largest experiments markets have ever handled. Prices may fall but we'll get clarity. We'll see which nodes were unavoidable.

Part of me thinks markets aren't able to handle this experiment in their current form. Maybe handling an experiment this important and pervasive requires evolving our markets themselves.

If bubbles are how we compute possible futures, the work is to choose which future we believe in. After margin finds its level and the hype passes, what do we want to survive?

In a present where every keystroke passes through a chain of factories and servers. The future will settle with the chain of the most unavoidable nodes. The ones you can't go around. They might be quiet now, but sometime sooner rather than later we'll notice our lives depend on them.

I think maybe a shortlist for those nodes looks something like this:

\begin{itemize}
    \item Where physics enforce adoption.
    \item Where taste enforces loyalty.
    \item Where memory compounds usefulness.
    \item Where governments guarantee demand.
    \item Where switching costs rise.
    \item Where the asset is still worth owning if the hype goes quiet.
\end{itemize}

\end{document}
